\chapter[Sermon 96. He Continues Yet in Describing the Blessed Seats.]{He Continues Yet in Describing the Blessed Seats.}
\label{sermon:096}
\markright{Sermon 96. He Continues Yet in Describing the Blessed Seats.}

And he showed me a pure river of water of life, clear as crystal, proceeding out of the seat of God and of the Lamb. In the midst of the street of it, and on either side of the river, was wood of life, bearing twelve kinds of fruit, and yielding fruit every month. The leaves of the wood served to heal the people. And there shall be no more curse, but the seat of God and the Lamb shall be in it, and his servants shall serve him. They shall see his face, and his name shall be on their foreheads. And there shall be no night there; they need no lamp, nor light of the sun, for the Lord God gives them light, and they shall reign forevermore.

\index{Augustine, Saint}\index{Zechariah (Book of)}\index{Ezekiel (Book of)}In the twelfth place is described John’s vision of the pleasantness, abundance, and richness of food in the City of God. Rivers make cities pleasant and delightful. Without springs, fountains, and wholesome waters, cities decay and are scarcely worthy of the name. But if they lack sustenance, they are utterly lost. Therefore, our heavenly City excels in all these things and is most noble. It not only provides sustenance but gives it with great pleasure and delightful finesse. For the trees in this City not only bear fruit but also give an unspeakable and inestimable joy. The river moreover runs through the midst of the streets; on the banks of either side are trees most beautiful to behold, bearing the fruits of life. As I have often intimated in this description, I repeat now that these things are not to be understood literally, as the Millenarians take them. For the Lord speaks to us, even “lisps,” so that we, with our limited understanding, might grasp these things. If anyone should desire earthly things, I believe he could desire nothing greater than what is described here. We should therefore consider that if the Lord could give these earthly things, and would, why could he not give greater things to the souls of the godly and bodies glorified? Yes, the Lord desires that, being withdrawn from the contemplation of earthly things, we should look altogether for celestial and divine things, worthy of blessed souls and glorified bodies. Which, indeed, how great they are and what they shall be, no human tongue can express, however eloquent. For the Lord has prepared greater things for his servants than we can comprehend. Therefore, he brings forth here ample material so that, in a certain manner, we might conceive of heavenly things much more excellent than they are. Therefore, the sense and meaning of all those things spoken here of the river of life and the wood of life, by a most excellent amplification, is none other than that the blessed in the heavenly country shall be quickened by God and preserved in that happy life with unending delight. And there is no doubt that John has borrowed these things, as he does all the rest, since he is the interpreter of the Prophets, from the Scriptures. Therefore, he alludes to Paradise, whose description is set forth in the second chapter of Genesis, which agrees very well with this description of Heaven. For there also springs a river in Paradise, which immediately is divided into four heads and waters the garden of pleasure most pleasantly. In that same Paradise is the wood, that is, the tree of life, bringing forth lively fruit to the eaters, as explained by Augustine in book 13 of \textit{The City of God}, chapter 20. But for the sin of our first parent, we were cast out of that Paradise; and Christ came to bring us again into Paradise, that is to say, into high felicity. Now, therefore, that true paradise, prepared for us by Christ, is shown in heaven and is here described. Into this Paradise entered the Lord after death, and brought with him into the same also the faithful thief, to whom he said, “Verily I say unto thee, this day shalt thou be with me in paradise.” Therefore, we ought not here to imagine the gardens of Hesperides in Earth, or in the air above the globe of the Moon, and reason of Paradise terrestrial. Our Paradise is celestial, which is prepared for us in heaven, as Paul has said in Philippians 3. Paradise is called a garden of pleasure, as at present it is called a golden City or of precious stones, verily by a trope on either side. Hereunto appertains also a passage from Zechariah in chapter 14. There is also another passage of Scripture in chapter 47 of Ezekiel, which John has translated or written out into this place in a manner word for word: “By the river,” he says, “on either side of it shall grow up all manner of trees that bear fruit, whose leaves shall not fall, nor the fruits fail, but every month shall they bring forth new fruit. For their waters flow out of the sanctuary, and their fruits shall be meat, and their leaves medicinable.” And Ezekiel, under a figure, sees that same blessed life and happy seats which John at present sees, by the showing of the Angel. And either of them both sees the happy seats after the same sort, and under the like figure. For there is one only blessedness, common to all the faithful of the whole world. The Patriarchs, Prophets, Apostles, and Martyrs achieve all one felicity. They see the River on either side, and the same flowing out of the Sanctuary, or seat of God. They see on either side the river, trees planted that bring forth the fruits of life. They bring forth fruit every month, fresh and new; and the leaves of either do heal. I suppose the old poets borrowed from the Scriptures such things as they wrote in verses concerning Ambrosia and Nectar, the meat and drink of the Gods. That short verse of Martial is known.

\index{John, Saint}And Grammarians derive those terms of immortality. But our Saint John here reasoning more elegantly and better concerning these matters, says how the angel showed him a river, which he commends to him by such properties as water is wont to be commended: for its purity, brightness, and clearness. He adds a parable, which sheds light on what he has said, and says, clear as a crystal. After this, he adds that this river is the river of the water of life, that is, lively water, which preserves those who drink it in life. Finally, he shows also the origin or springhead of this river, deriving it from the seat of God, of which seat or throne I have spoken in the fourth and fifth chapters of this book. And by all these things is signified nothing else but that life proceeds from God alone, which he gives to those who serve him in that blessed country, pure, clear, bright, most tested and most perfect, and altogether divine. Concerning the lively springs and fountains of waters, we have touched somewhat in the end of the seventh chapter of this book. Note again that God and the Lamb are so joined together that no man unless he be mad will deny the Son to be of the same substance with the Father.

Now follows the description of this divine city. The sustenance in the eternal land is the tree of life. It is the Hebrew phrase, “the wood of life,” for the tree of life, or lively sustenance. For there is added bearing fruit. Whether you understand that John saw only one tree, as also there was only one tree in paradise: or more, as in Ezekiel, so that by the general word we may understand the particular kinds of trees, it shall be all one. The location of the tree he shows diligently, to be set in the midst of the street of the city, and on either side the river (whereby it is doubtless gathered that there were many trees) to wit, on the banks of the river, that they may draw lively nourishment out of the river, which flows from the throne. And hereby I suppose is signified that the heavenly food is common and free for all, not locked up, or kept for a few. It is found in the midst of the street of the city: thus the sustenance stands open, and is not hidden. And it draws a lively force out of the river, which springs out of the seat. For that heavenly life is from God, and flows to all his elect. Moreover, it is also declared most diligently, what manner of fruit this shall be: the tree of life says he, does bear fruit twelve times in the year, so that every month it bears fruit fresh and new. The first fruits are delightful: and those who commonly abhor old fruit, had rather have new. Therefore in that blessed land there shall be nothing tedious, unpleasant, loathsome, or in any wise to be rejected, but all things shall be most pleasant, most delicate, fresh, and delectable.

Now he does not overlook the leaves, and, as Ezekiel did, demonstrates their use. He says they serve as a medicine to heal the people. This is not because there will be diseases or sores in that blessed land, but so that the blessed may have continual and perpetual health. He calls these people Gentiles, not that the Gentiles are presently unclean, but because they once were, but now, being purged by Christ, live whole and sound forevermore.

\index{Antichrist}And by these allegories he has hitherto depicted in parts those blessed dwellings, prepared for the faithful in that everlasting realm, under the image of a most noble city. Which after he has shown us, he seems as it were to have opened heaven itself, and set forth the eternal felicity to be seen in a manner with mortal eyes, and even to have pointed with the finger, to no other end than that we should be strong and steadfast in the faith of our Lord Jesus Christ. And should never think that one who has ever seen those blessed dwellings, to which we are called by the denying of all pleasures? What if you should despise the present pleasures and should obtain none in time to come? This thought is wicked. Faith teaches you otherwise. But what say you more? Do you desire to know and see such things as God has shown you? You have seen enough and abundantly at this present. The Lord has shown you abundantly enough of celestial life and pleasure at this present. Endeavor now only, that the devil, the world, and Antichrist trodden under, you may aspire and be lifted up into those heavenly dwellings. Moreover, beware that you be not more curious than is fitting or required, and that you seek not to know more, and more exact things of the heavenly tower and perpetual joys, than the Lord himself, who alone knows these things, has revealed to you. Let this evident demonstration of eternal life suffice us. I believe never anyone has disputed better or more rightly, more elegantly and more evidently of the blessed life than here Saint John has done. Let us therefore repose ourselves in God, let us believe his words, let his revelation suffice us, and let us desire to be joined with him in this heavenly court, in all felicity and eternal life most perfect.

And now John, recalling the main points of this matter, and concluding this passage concerning eternal life, finishes this everlasting felicity in seven parts. We will touch on these briefly, for many believe we have already spoken of them sufficiently and abundantly. First, there will be no curse, no condemnation, no malediction, neither war, nor famine, nor disease, nor anything as is recited by Moses among the curses in Deuteronomy 27:28. It is not that all who are subject to these things are accursed. Job and other holy men were afflicted with sickness; but commonly, the unbelievers and wicked are plagued with them. Not that they should be disciplined and profit in godliness, but that they should first be afflicted here, and by certain degrees pass on to greater torments. What then?

The second part adds: but the throne of God and of the Lamb will be in that city. That is, the kingdom of God will be there, and God will reign, and all blessing, no malediction, will be among the elect. Therefore, whatever joyful things the Prophets, Christ, and the Apostles have spoken of the kingdom of God, those things will be in heaven, and the blessed will experience them. And again, the Father and the Son are joined together inseparably in the unity of essence, which nevertheless, in the distinction of persons, are exceedingly well not divided, but discerned. These mysteries of the blessed Trinity are known to the faithful.

Here follows the third point. Some may wonder what the blessed will do in the world everlasting. Therefore John and his associates say that they shall serve God, whom I say is the Lamb, and they shall worship him in honoring, praising, and magnifying him forever. Therefore they shall wholly devote themselves to godly worship. This thing will indeed be to him great pleasure, as Augustine also shows in another place.

\index{Tertullian}Fourthly, they shall see the face of God. Augustine discusses much of what is said about God to Paulina, and warns the faithful not to imagine carnal things here. Moses in Exodus 33 and John the Apostle in John 14 have considered it the greatest blessing to see God as he is, and as is commonly said, face to face. And there is undoubtedly in this sight and enjoyment, great felicity, everlasting and most complete joy; however, in this present world, as the Lord said to Moses, it happens to no one. The holy fathers have indeed seen God, but by a form, and so far as he has deigned to reveal and show himself to be seen. As Tertullian shows in the book \textit{Against Praxeas}, but to see the full glory of God with inestimable joy will be granted to us first when, being delivered from misery and purged from corruption, we shall also be clarified in body; then, as John also said in 1 John 3, we shall see him as he is. Job, the most righteous man, speaking of this vision of God, said: when they shall have encircled or clothed (that is, the Father, Son, and Holy Spirit) this (namely, my body), I shall behold God outside of my flesh; whom I shall see for myself, and my eyes shall look upon, and no other. This is my only desire. Paul, the Doctor of the Gentiles, also spoke of this saying, and said: Now we see in a mirror, even in a dark speaking; but then shall we see face to face, \&c. And Augustine has also discussed this vision in his book \textit{The City of God}, toward the end.

\index{Primasius}Finally, they shall have the name of God in their foreheads, either because they will be the children of God, as we have heard in the Epistle to the Philadelphians, in chapter 3 of this book. Indeed, in the heavenly realm it will be plainly known to all who are the children of God. In this world, they are commonly regarded as the children of the devil, yet in truth are the children of God. But this will clearly appear in another world, to the great glory of the elect. And indeed, the brightness of God will shine from the foreheads, or countenances, of the elect, as in times past the brightness of the Lord shone from the face of Moses and Christ. Or because all saints will know one another, since the virtue of God rests in their countenances—a sense I perceive has pleased Primasius. Or because they will be priests before the Lord forever, as the prophets have taught of the chosen. In ancient times, the high priest bore the very name of God on his forehead on a plate of gold, bound to his head with a lace. Undoubtedly, in the heavenly realm, the glory of the children of God will be wonderfully great, especially among those who have confessed the name of Christ on earth; for these, the heavenly Father will glorify.

In the sixteenth member, what has been said once or twice before is repeated: that the elect in heaven are illuminated with divine glory, of which enough has been spoken previously.

In the last and seventh part, encompassing as it were all things of life and felicity, and uttering with one word: they shall reign, he says, forevermore. The Lord Jesus grant to us his faithful, that such things as we have now heard plentifully from his mouth, we may shortly experience in our souls and bodies, and may cry with joy, to God the Father most merciful, and to Jesus Christ the Redeemer most mighty and benign, and to the Holy Spirit the most sweet Comforter be praise and glory forevermore. Amen.
